\chapter{L'ingresso generale: una marcia per uno scarico non
uniforme}
\label{ch:generalinlet}

Il capitolo precedente si è chiuso su una frontiera: lo swirl a
vortice libero conserva esattamente la struttura caratteristica, e
qualunque cosa oltre esso rende il flusso \emph{rotazionale}, che la
marcia certificata non sapeva rappresentare. Questo capitolo rimuove
quella frontiera per il caso che conta per un RDE, e lo fa portando
qui una macchina che il codice arbitro possiede da sempre.

\section{Che cosa la marcia certificata assumeva, e perché contava}

Ogni nodo della marcia, finora, conservava quattro numeri: la
posizione e le due componenti di velocità, $(x, y, u, v)$. Nulla di
termodinamico veniva conservato, perché nulla serviva: lo stato del
gas era recuperato dalla \emph{sola velocità},
\begin{equation}
  h_t = h_0 - \tfrac12 q^2 \;\longrightarrow\; T
  \;\longrightarrow\;
  p = p_{\mathrm{ref}}\exp\!\Big(\frac{s(T) - s_0}{R}\Big),
  \label{eq:homentropic}
\end{equation}
con una singola entalpia di ristagno $h_0$ e una singola entropia
$s_0$ per l'intero campo. È l'ipotesi \emph{omoentropica}: un solo
serbatoio, una sola isentropica, ogni linea di corrente attinta dallo
stesso gas nello stesso stato.

Per un ugello alimentato da una camera di calma questo è esattamente
corretto. Per un motore a detonazione rotante è una scelta di
modello. L'anello reale consegna uno scarico che non è né uniforme né
privo di urti: la detonazione passa, i prodotti caldi seguono quelli
più freddi, e ogni linea di corrente arriva con la propria
temperatura totale e la propria entropia. La marcia non poteva
accettare dati simili --- non perché i numeri fossero difficili da
fornire, ma perché sarebbero stati \emph{sovradeterminati}. Dato un
profilo $(p, T, M)$ riga per riga, la marcia può onorare la velocità
e la direzione; la pressione che poi predice dalla
\eqref{eq:homentropic} è quella che dice l'isentropica globale, e se
la $p$ fornita non concorda non c'è dove mettere la differenza. Due
numeri su tre entrano; il terzo non ha casa.

\section{Che cosa l'arbitro faceva già}

Il codice Fortran indipendente ha varcato questa frontiera molto
tempo fa. Il suo modulo generico delle caratteristiche risolve il
classico processo elementare \emph{rotazionale}: entropia ed entalpia
di ristagno sono trasportate come \textbf{invarianti di linea di
corrente}, ogni cella localizza la linea di corrente che passa per il
suo nuovo punto, e lo stato al piede di quella linea è ricostruito da
\emph{invarianti interpolati} anziché da pressioni e densità
interpolate --- una distinzione che il registro di sviluppo di quel
codice annota come correzione, perché interpolare le grandezze
primitive fa uscire dalla varietà termodinamica mentre interpolare
gli invarianti non può farlo.

La conseguenza è esattamente la capacità che ci mancava: l'arbitro
legge da file un'intera linea iniziale punto per punto --- posizione,
densità, velocità, angolo, pressione, temperatura, Mach, $\gamma$ ---
e la marcia. Non c'era ragione di inventare una seconda risposta a un
problema già risolto. Questo capitolo porta qui la loro.

\section{Le tre modifiche}

\paragraph{I nodi portano la propria termodinamica.} Un nodo diventa
$(x, y, u, v, s, h_0)$: i due invarianti si uniscono ai quattro numeri
cinematici. La chiusura~\eqref{eq:homentropic} è invariata nella
forma e generalizzata nel contenuto --- ogni nodo usa i \emph{propri}
$(s, h_0)$ là dove la chiusura certificata usava la coppia globale.
Alla coppia globale le due coincidono identicamente, e il carrier
misura quell'identità anziché assumerla: la deviazione è esattamente
nulla, perché il percorso aritmetico è lo stesso.

\paragraph{Le equazioni di compatibilità cambiano forma.} Le celle
certificate risolvono la relazione di compatibilità gasdinamica
scritta in $(u, v)$. Quella forma è derivata assumendo flusso
irrotazionale, e per il teorema di Crocco i gradienti di entropia
\emph{sono} vorticità, quindi non sopravvive alla stratificazione. Il
processo elementare rotazionale usa invece la relazione lungo le
linee di Mach in pressione e angolo del flusso,
\begin{equation}
  \mathrm{d}\theta \pm Q\,\mathrm{d}p + S\,\mathrm{d}x = 0,
  \qquad
  Q = \frac{\sqrt{M^2-1}}{\rho q^2},
  \qquad
  S = \frac{\delta \sin\theta}{y\,M\,\cos(\theta \mp \alpha)},
  \label{eq:ptheta}
\end{equation}
che non richiede irrotazionalità: $Q$ è il coefficiente acustico
lungo una linea di Mach e $S$ è la sorgente assialsimmetrica che si
annulla nel flusso piano. Ogni coefficiente è valutato nello stato al
punto medio della propria corda, la forma implicita che le celle
certificate già usano.

\paragraph{La linea di corrente entra fra le incognite.} Una cella
interna risolve ora cinque numeri, $(x_4, y_4, u_4, v_4, t)$: la
posizione dai due incroci di linee di Mach, lo stato dalle due
relazioni~\eqref{eq:ptheta}, e $t$ --- il punto in cui la linea di
corrente che arriva al nuovo nodo attraversa la linea di dati a
monte. Gli invarianti nel nuovo nodo sono quelli della linea di dati
interpolati a quel $t$. Risolvere il piede \emph{insieme} al punto,
anziché iterare attorno a esso, è l'espressione naturale della stessa
fisica nell'ambiente a celle implicite di questo programma.

La parete e la frontiera libera non richiedono equazioni nuove:
entrambe sono linee di corrente, quindi ciascuna porta con sé, per
l'intera marcia, gli invarianti della riga da cui parte.

Con tutto ciò a posto, una linea $(p, T, M, \theta)$ riga per riga è
\textbf{esattamente determinata} --- tre numeri termodinamici per i
tre gradi di libertà dello stato $(s, h_0, q)$, più la direzione ---
e ripercorre la chiusura in andata e ritorno a quindici cifre.

\section{Due difetti colti dai test}

Entrambi sono registrati perché entrambi producevano uscite
plausibili.

\textbf{Il piede della linea di corrente usciva dalla propria corda.}
La prima versione interpolava gli invarianti lungo la corda che
congiunge i due punti di dati della cella. Su un getto stratificato la
marcia si rifiutava allora di convergere: raffinando la linea
iniziale l'errore restava dov'era, vicino a una parte su mille. La
sonda cella per cella ha mostrato perché. Ogni colonna scende di
circa $\tan\mu$ per la spaziatura delle stazioni, quindi appena le
righe sono abbastanza fini perché questo scorrimento superi la
spaziatura fra righe, la linea di corrente che arriva a un nodo cade
\emph{fuori} dalla corda fra i suoi due vicini --- e interpolare a
$t<0$ è estrapolare. Gli invarianti venivano allora trasportati lungo
una corda che non li conteneva. L'arbitro non ha questo difetto: la
sua cella cerca l'intersezione vera sull'intera colonna a monte.
Portata qui --- come \emph{decisione registrata}, nello stesso modo
in cui il piede di parete è sempre stato registrato --- l'errore di
trasporto nel corpo del getto è calato di due ordini di grandezza ed
è comparsa la convergenza (Figura~\ref{fig:s24rot}c).

\textbf{Un rigettatore va posto dove le equazioni portano carico.} La
verifica che corrompe un segno di compatibilità è stata dapprima
eseguita sul getto parallelo esatto, e \emph{ha promosso il codice
corrotto}: su quel mondo le differenze di pressione si annullano
identicamente, quindi il termine corrotto moltiplica zero e la marcia
corrotta riproduce la risposta giusta. Un rigettatore che non può
fallire non dimostra nulla. Spostato sul mondo del plug, dove le
stesse equazioni fanno lavoro reale, separa di tre ordini di
grandezza.

\section{Che cosa stabiliscono i test}

\begin{figure}[htbp]
\centering
\includegraphics[width=\textwidth]{fig_s24_rotational.pdf}
\caption{L'ingresso generale, dall'esecuzione di riferimento. (a)~Una
linea iniziale stratificata: entropia, entalpia di ristagno,
temperatura e numero di Mach variano tutti attraverso il getto,
ciascuna curva normalizzata sul proprio intervallo con l'intervallo
stampato --- il tipo di dati che la marcia certificata non poteva
accettare. (b)~Il campo di entropia che la marcia trasporta
attraverso una spina: gli strati cavalcano le linee di corrente dalla
linea iniziale all'uscita, ed è questo l'aspetto che ha ``l'entropia è
un invariante di linea di corrente'' quando un solutore lo rispetta.
(c)~Verifica: l'errore di trasporto rispetto alla soluzione
stratificata esatta cala con il numero di righe della linea iniziale
all'ordine dello schema, mentre l'artefatto, riportato a parte, della
regola di crescita delle righe della maglia su un getto che non si
allarga non cala --- è topologia, non trasporto, ed è confinato a una
banda della colonna che la verifica esclude e dichiara.}
\label{fig:s24rot}
\end{figure}

La verifica usa una soluzione esatta per delimitare l'errore. Un
getto rettilineo a pressione uniforme, con temperatura e numero di
Mach diversi su ogni linea di corrente e senza angolo di flusso, è
una soluzione esatta delle equazioni assialsimmetriche: nulla ruota,
nulla si espande, e ciascuna linea di corrente si limita a portare a
valle il proprio stato. Se gli invarianti variano \emph{linearmente}
attraverso il getto anche l'interpolazione di cella è esatta, quindi
la marcia deve riprodurre la risposta all'arrotondamento --- misurato
$3\times10^{-8}$ su cinquecento celle certificate. Curvando la
stratificazione l'interpolazione diventa del secondo ordine, e
l'errore cala allora al raffinarsi delle righe:
$2.2\times10^{-5} \to 6.5\times10^{-6} \to 7.5\times10^{-7}$ su
$21 \to 41 \to 81$ righe, più rapidamente del secondo ordine che lo
schema promette.

Sul mondo del plug, a invarianti uniformi, le celle rotazionali e
quelle certificate sono due discretizzazioni delle stesse equazioni
continue, quindi non devono coincidere esattamente --- e
l'affermazione onesta è che la loro differenza \emph{si riduce}
all'ordine comune dei due schemi, da $3.0\times10^{-5}$ a
$8.8\times10^{-6}$ sotto raffinamento congiunto. Su una spina
genuinamente stratificata la marcia certifica a $0.058$, conserva la
massa a $8\times10^{-3}$ dopo la prima colonna, e trasporta il flusso
di entropia con una deriva di $5\times10^{-4}$ --- l'invariante
trasportato, misurato anziché assunto.

Un artefatto è riportato anziché assorbito. Su un getto che
\emph{non} si allarga, la regola di maglia che aggiunge una riga per
colonna --- scritta per uno scarico che si espande --- ammassa le
vecchie righe di bordo contro la frontiera stazionaria e apre un
vuoto in mezzo alla colonna. Le corde che attraversano quel vuoto
interpolano un profilo curvo lungo mezzo metro, e l'errore che ne
risulta non cala con il raffinamento perché è fissato dalla
topologia, non dalla spaziatura. È confinato alla banda ammassata, è
stabile all'uno percento sull'intera scala di raffinamento, e
l'affermazione sul trasporto è valutata al di sotto di esso. Le
geometrie di plug che questo programma progetta davvero sono scarichi
che si allargano e non lo sviluppano mai.

\section{Che cosa apre}

La conseguenza immediata è che uno scarico di RDE con vene calde e
fredde può essere marciato così com'è, anziché nella sua caricatura
uniforme in media; e poiché la linea iniziale che il port accetta
porta lo stesso contenuto del file dei valori iniziali dell'arbitro,
uno stesso ingresso può ora alimentare entrambi i codici --- il che
rende un confronto fra codici su un ingresso \emph{rotazionale} il
prossimo test naturale. Restano aperte e non misurate due ulteriori
questioni: se l'aggiunto del Capitolo~\ref{ch:spline} differenzi
attraverso una marcia stratificata così com'è (le celle passano per
la stessa macchina a solutore implicito, quindi dovrebbe), e se swirl
e stratificazione si compongano --- il vortice libero richiede
entalpia totale uniforme per il proprio argomento di
irrotazionalità, e uno scarico stratificato con swirl è
un'affermazione diversa, e più difficile.
